
В рамках данной работы была решена задача разработки и сравнительного анализа четырёх алгоритмов автоматического построения расписания учебных занятий, основанных на принципиально разных подходах к оптимизации: генетическом алгоритме, имитации отжига, поиске с запретами и точном программировании ограничений. Все четыре алгоритма были реализованы на единой модели данных, описывающей дисциплины, преподавателей, аудитории и группы, и оценивались единым критерием качества решения, различающим жёсткие ограничения, нарушение которых делает расписание физически нереализуемым, и мягкие ограничения, отражающие желательные, но не обязательные свойства расписания.

Поставленные во введении задачи были выполнены последовательно и в полном объёме. Изучение теоретических основ задачи составления расписания и существующих подходов к её решению позволило обоснованно выбрать четыре рассмотренных алгоритма как представителей принципиально разных алгоритмических парадигм — от точного метода, гарантирующего оптимальность или доказывающего отсутствие решения, до метаэвристик, жертвующих этой гарантией в пользу гибкости и масштабируемости. Единая модель данных и единый критерий оценки качества, спроектированные до начала реализации конкретных алгоритмов, оказались необходимым условием честного сравнения: без них результаты разных алгоритмов были бы выражены в несопоставимых единицах измерения. Реализация всех четырёх алгоритмов сопровождалась рядом содержательных трудностей, часть которых была обнаружена не на этапе проектирования, а эмпирически — в частности, дефект, связанный с ограничением размера итогового расписания, проявившийся одинаково в трёх из четырёх реализаций и обнаруженный благодаря логическому противоречию между ответом точного решателя и ответами эвристических методов на одном и том же наборе данных. Этот эпизод сам по себе стал практическим подтверждением одного из выводов работы: наличие в распоряжении точного метода ценно не только как самостоятельный инструмент решения задачи, но и как средство проверки корректности эвристических реализаций.

Экспериментальное сравнение на наборах данных разного размера и структурной сложности показало, что поиск с запретами обеспечивает наилучшее соотношение качества решения и затраченного вычислительного времени среди трёх эвристических методов, имитация отжига является самым быстрым методом при умеренном качестве решения, а генетический алгоритм в реализованном виде уступает обоим по качеству при существенно более высоких вычислительных затратах. Решатель ограничений на разрешимых наборах данных находил решение, не худшее, чем любой из эвристических методов, а на структурно неразрешимом наборе данных — единственный из четырёх алгоритмов дал не приближённый, а доказанный ответ.

Практическая значимость работы заключается в том, что предложенная методика сравнения — единая модель данных, единый критерий оценки качества, проверка на нескольких независимых осях (случайность запуска и структура исходных данных) — может быть использована как основа для выбора конкретного алгоритма при построении полноценной системы автоматизации расписания, в зависимости от ожидаемого размера задачи и доступного времени на вычисление. Для образовательного учреждения с умеренным числом дисциплин и групп предложенные результаты прямо указывают на целесообразность использования точного метода либо, при больших объёмах данных, поиска с запретами как наиболее сбалансированной альтернативы.

Направления дальнейшего развития работы включают устранение выявленной, но не успевшей проявиться в эксперименте архитектурной неточности табу-списка в поиске с запретами, переход от взвешенной суммы мягких критериев к строго лексикографической многоступенчатой оптимизации в решателе ограничений там, где это содержательно оправдано, а также расширение модели данных поддержкой нескольких групп с пересекающимися ресурсами — преподавателями и аудиториями, — что превратило бы задачу из однопоточной в полноценную многогруппную постановку, ближе соответствующую масштабу реального учебного заведения.
